% =========================================================================
% CAPÍTULO 11
% =========================================================================
\chapter{Conclusiones y Recomendaciones}

Se cumplió el objetivo general: se diseñó, implementó y validó empíricamente un sistema de alerta temprana híbrido. El modelo XGBoost fue el de mejor desempeño en regresión ($R^2=0.6272$); el ensamble Stacking con SMOTE-ENN fue efectivo para la clasificación de insolvencia (ROC-AUC=0.9227). El aprendizaje automático aportó capacidad de generalización transversal entre instituciones; la econometría clásica aportó un punto de referencia riguroso para instituciones con historia extensa. Un sistema de alerta temprana basado en estos componentes es viable como herramienta complementaria —no sustitutiva— de la supervisión prudencial tradicional.

\section{Comprobación de Hipótesis}
Con base en los resultados del Capítulo 6, \textit{no se rechaza H1}: los modelos panel multivariados demuestran capacidad predictiva y de generalización superior a los modelos univariados en el sentido de cobertura institucional, aunque el modelo univariado ARIMA presenta menor error puntual en el subconjunto reducido de instituciones donde es aplicable. Respecto a H1a, los resultados \textit{no ofrecen evidencia fuerte de soporte}: las variables macroeconómicas exógenas contribuyeron de forma secundaria, aunque no nula, frente al peso dominante de las variables contables rezagadas, para el horizonte de un trimestre evaluado. Este hallazgo es consistente con la economía financiera: a horizontes cortos, la inercia contable domina; a horizontes más largos, el entorno macroeconómico tendría mayor relevancia.

\section{Recomendaciones}
\textbf{Académicas:} extender el análisis a horizontes múltiples y reemplazar la partición por posición por una partición cronológica por periodo. \textbf{Técnicas:} investigar la anomalía de NICAP con la fuente de datos original antes de cualquier uso regulatorio del modelo. \textbf{Para SOFIPOs:} utilizar el sistema como insumo de alerta interna, no como sustituto del juicio experto de comités de riesgo. \textbf{Para la CNBV:} explorar la integración piloto del pipeline como herramienta de apoyo a la supervisión basada en riesgo. \textbf{Para investigaciones futuras:} ampliar el panel conforme se acumulen nuevos trimestres regulatorios y evaluar arquitecturas LSTM cuando la profundidad temporal lo permita.

